\documentclass[framenumber,t]{beamer}
\usepackage[absolute,overlay]{textpos}
\usepackage[T1]{fontenc}
\usepackage[utf8]{inputenc}
\usepackage{booktabs}
\usepackage{multicol}

\usepackage{amsmath,amsfonts,amssymb}
\usepackage{bm}

%----------------------------
% fonts
\usepackage{helvet}
\usepackage{mathptmx} % fatter math font
\usefonttheme[onlymath]{serif}
\usepackage{lmodern}
\renewcommand{\ttdefault}{cmtt}

%----------------------------
% colors
\definecolor{fhnwblue}{RGB}{0,51,121}
\definecolor{orange}{RGB}{199,90,9}
\definecolor{orange_dark}{RGB}{162,74,14}

\setbeamercolor{frametitle}{fg=fhnwblue}

\setbeamercolor{section in toc}{fg=fhnwblue,bg=white}
\setbeamercolor{item projected}{bg=fhnwblue}
\setbeamertemplate{enumerate items}[default]
\setbeamertemplate{navigation symbols}{}
\setbeamercovered{transparent}
\setbeamercolor{enumerate item}{fg=fhnwblue}
\setbeamercolor{enumerate subitem}{fg=fhnwblue}
\setbeamercolor{enumerate subsubitem}{fg=fhnwblue}
\setbeamercolor{itemize item}{fg=fhnwblue}
\setbeamercolor{itemize subitem}{fg=fhnwblue}
\setbeamercolor{itemize subsubitem}{fg=fhnwblue}

\setbeamerfont{framesubtitle}{size=\large}

%----------------------------
% alerted text
\setbeamerfont{alerted text}{series=\bfseries}
\setbeamercolor{alerted text}{fg=orange_dark}

%----------------------------
% code command (can deal with '$', '_', etc.)
\makeatletter
\newcommand\code{\bgroup\@makeother\_\@makeother\~\@makeother\$\@makeother\^\@codex}
\def\@codex#1{{\normalfont\ttfamily\hyphenchar\font=-1 #1}\egroup}
\makeatother

% proglang
\let\proglang=\textsf

%----------------------------
% itemize without identation
\setlength{\leftmargini}{7pt}
\setbeamertemplate{itemize items}{\scalebox{0.6}{$\blacksquare$}}

%----------------------------
% code listing (needs specification 'fragile')
\usepackage{listings}
\lstdefinestyle{in_big}{
    language=R,
    basicstyle=\small\ttfamily\bfseries\color{orange_dark}
}
\lstdefinestyle{in_small}{
    language=R,
    basicstyle=\footnotesize\ttfamily\bfseries\color{orange_dark}
}
\lstdefinestyle{out_big}{
    language=R,
    basicstyle=\small\ttfamily\color{fhnwblue}
}
\lstdefinestyle{out_small}{
    language=R,
    basicstyle=\footnotesize\ttfamily\color{fhnwblue}
}

%----------------------------
\usepackage{tikz}
\usetikzlibrary{backgrounds, positioning}
\usepackage{pgf}
\logo{\pgfputat{\pgfxy(-6.2,8.4)}{\pgfbox[right,base]{%
    \includegraphics[height=0.65cm]{FHNW_HW_E_10mm.pdf}}}}
\newcommand{\nologo}{\setbeamertemplate{logo}{}}
\makeatletter
\setbeamertemplate{frametitle}{
    \bfseries
    \ifbeamercolorempty[bg]{frametitle}{}{\nointerlineskip}%
    \@tempdima=\textwidth%
    \advance\@tempdima by\beamer@leftmargin%
    \advance\@tempdima by\beamer@rightmargin%
    \hspace*{1.25em} %%%%%%%%%%%%% For example insert shift to right
    \begin{beamercolorbox}[left,wd=\the\@tempdima]{frametitle}
        \usebeamerfont{frametitle}%
        \vbox{}\vskip 4.5ex%
        \if@tempswa\else\csname beamer@ftecenter\endcsname\fi%
        \strut\insertframetitle\strut\par%
        {%
            \ifx\insertframesubtitle\@empty%
            \else%
            {\usebeamerfont{framesubtitle}\usebeamercolor[fg]{framesubtitle}%
                \insertframesubtitle\strut\par}%
            \fi
        }%
        \vskip-1ex%
        \if@tempswa\else\vskip-.3cm\fi% set inside beamercolorbox...
        \vspace{0.5em}
    \end{beamercolorbox}%
}
\makeatother

\setbeamertemplate{footline}[text line]{%
    \parbox{\linewidth}{%
        \vspace*{-15pt}%
        \insertshorttitle\hfill\insertshortauthor\hfill\insertpagenumber}
}
\setbeamertemplate{navigation symbols}{}

% -----------------------------------------------------------------------------
\begin{document}

\title[robsurvey: useR! 2022 conference]{\emph{robsurvey:} Robust Survey Statistics Estimation}
\subtitle{useR! 2022 Conference}

\author[Tobias Schoch]{Tobias Schoch (joint work with Beat Hulliger)}
\date{\today}

% -----------------------------------------------------------------------------
% title page
{
\setbeamertemplate{footline}{}
\begin{frame}
   \pgfputat{\pgfxy(-1,-7.6)}{\pgfbox[left,base]{%
        \includegraphics[width=12cm]{title.png}}}
    %
    \begin{textblock}{20}(1.15,3.6)
        \Large\bfseries\color{fhnwblue}\inserttitle
    \end{textblock}
    %
    \begin{textblock}{20}(1.15,4.8)
        \small\bfseries\color{fhnwblue}\insertsubtitle
    \end{textblock}
    %
    \begin{textblock}{20}(1.15,15.2)
        \scriptsize\color{black}\insertauthor, ~\insertdate
    \end{textblock}
    %
    \vspace*{3.3cm}\hspace{-1.2cm}%
    \begin{tikzpicture}[ scale=1]
    \draw [fhnwblue,fill] (0,0) rectangle (1.1,4);
    \end{tikzpicture}
\end{frame}
}
\addtocounter{framenumber}{-1}
% -----------------------------------------------------------------------------
\begin{frame}[fragile]
\frametitle{Outline}
    \begin{itemize}
        \setlength\itemsep{0.75em}
        \item Basic robust estimators
        \item Robust weighted regression
        \item Summary \& outlook
    \end{itemize}
    \vspace{3em}
    \begin{minipage}{6cm}
        \alert{Requirement}
    \end{minipage}
    \begin{lstlisting}[style=in_big]
> library("robsurvey", quietly = TRUE)
> data("losdata")
> data("counties")
    \end{lstlisting}
\end{frame}

% -----------------------------------------------------------------------------
\bgroup
\setbeamercolor{background canvas}{bg=fhnwblue}
\begin{frame}[c,plain]{}
\begin{center}
    \textcolor{white}{\large PART 1}\\
    \vspace{1em}
    \textcolor{white}{\Huge Basic Robust Estimators}
\end{center}
\end{frame}
\egroup

% -----------------------------------------------------------------------------
\begin{frame}[fragile]
\frametitle{The \code{losdata} data}
    \begin{itemize}
        \setlength\itemsep{0.5em}
        \item \alert{Length of stay} (LOS) in hospital (days per year)
        \item \alert{Sample} of $n = 71$ patients (population size $N = 2479$)
    \end{itemize}
    \begin{lstlisting}[style=in_big]
> head(losdata, 3)
    \end{lstlisting}
    \begin{lstlisting}[style=out_big]
  los   weight  fpc
1  10 34.91549 2479
2   7 34.91549 2479
3  21 34.91549 2479
    \end{lstlisting}
    where\\
    \vspace{0.5em}
    \begin{tabular}{ll}
        \code{los} & length of stay in hospital\\
        \code{weight} & sampling weight\\
        \code{fpc} & population size (finite population correction)\\
    \end{tabular}
\end{frame}

% -----------------------------------------------------------------------------
\begin{frame}
\frametitle{The \code{losdata} data {\small (cont'd)}}
    \begin{itemize}
        \setlength\itemsep{0.1em}
        \vspace{0.5em}
        \item \alert{Distribution} of \code{los}
        \begin{itemize}
            \setlength\itemsep{0.5em}
            \item Average: 13.03 (days)
            \item Maximum: 182
        \end{itemize}
    \end{itemize}

    % library(robsurvey)
    % attach(losdata)
    % file_name <- "boxplot_los.pdf"
    % pdf(file = file_name, height = 3)
    % boxplot(los, horizontal = TRUE, xlab = "LOS (days)", cex.axis = 1.2,
    %     cex.lab = 1.4)
    % dev.off()
    % system(paste("pdfcrop.exe", file_name, file_name))
    \begin{center}
        \includegraphics[height=3cm]{figure/boxplot_los}
    \end{center}
    % \begin{itemize}
    %     \setlength\itemsep{0.5em}
    %     \item $\mathrm{var}(\bar{y})$ is large
    %     \item Is there an \alert{estimator $\theta(y_1, \ldots, y_n)$} of the
    %         mean such that $\mathrm{mse}(\theta) < \mathrm{var}(\bar{y})$?
    %         E.g., winsorized mean or $M$-estimator?
    %     % \item Chambers (1986)
    %     %     \begin{itemize}
    %     %         \setlength\itemsep{0.5em}
    %     %         \item \alert{Representative outlier}: extreme but \alert{correct}
    %     %             value that \alert{represents} other population units
    %     %             similar in value
    %     %         \item \alert{Nonrepresentative outlier}: extreme observation
    %     %             whose value is either deemed erroneous or \alert{unique}
    %     %             in the sense that there is no other unit like it
    %     %     \end{itemize}
    % \end{itemize}
\end{frame}

% -----------------------------------------------------------------------------
\begin{frame}[fragile]
\frametitle{Two ``flavors'' of basic robust estimators\vspace{0.5em}}
\begin{itemize}
    \item \alert{bare-bone} functions: \code{weighted_mean} and
        \code{weighted_total}
    \item \alert{survey} methods: \code{svymean} and \code{svytotal}
\end{itemize}
\vspace{0.5em}
followed by (suffix)
\vspace{1em}
\begin{tabular}{ll}
    \addlinespace
    &\code{_winsorized()} and \code{_k_winsorized()}: winsorization\\
    \addlinespace
    &\code{_trimmed()}: trimming\\
    \addlinespace
    &\code{_huber()} and \code{_tukey()}: $M$-estimators\\
    \addlinespace
    &\code{_dalen()}: Dal{\'e}n's estimators: censoring of value and weight\\
\end{tabular}
E.g., \code{weighted_mean_winsorized()}
\end{frame}

% -----------------------------------------------------------------------------
\begin{frame}[fragile]
    \frametitle{Bare-bone functions}
    The 2\% one-sided \alert{winsorized weighted mean} is

    \begin{lstlisting}[style=in_big]
> attach(losdata)
> weighted_mean_winsorized(los, weight, LB = 0,
                           UB = 0.98)
    \end{lstlisting}
    \vspace{-1em}
    \begin{lstlisting}[style=out_big]
[1] 11.40845
    \end{lstlisting}
    \vspace{0.5em}
    \begin{itemize}
        \setlength\itemsep{0.5em}
        \item Lower bound \code{LB} = 0
        \item Upper bound \code{UB = 0.98} \alert{$\boldsymbol{\Rightarrow}$}
            largest 2\% of the observations are winsorized
        \item Return value: estimate (scalar)
        \item Light-weight, minimalistic, bare-bone
        \item Useful for package developers
    \end{itemize}
\end{frame}

% -----------------------------------------------------------------------------
\begin{frame}[fragile]
\frametitle{Survey methods}
\begin{lstlisting}[style=in_big]
> library("survey")
> dn <- svydesign(ids = ~1, fpc = ~fpc,
                  weights = ~weight,
                  data = losdata)
\end{lstlisting}
%where \code{survey::svydesign()} defines a survey design
\vspace{1em}
\begin{lstlisting}[style=in_big]
> svymean_winsorized(~los, dn, LB = 0, UB = 0.98)
\end{lstlisting}
\vspace{-0.5em}
\begin{lstlisting}[style=out_big]
     mean  SE
los 11.41 1.5
\end{lstlisting}
\vspace{0.5em}
\begin{itemize}
    \setlength\itemsep{0.5em}
    \item Computes \alert{standard errors (\code{SE})}
        using functionality of the \code{survey} package (Lumley, 2010, 2021)
    \item Return value: an instance of class \code{svystat_rob}
\end{itemize}
\end{frame}

% -----------------------------------------------------------------------------
\begin{frame}
\frametitle{Survey methods {\small (cont'd)}\vspace{0.25em}}
\framesubtitle{Utility methods}

\begin{tabular}{llr}
    \code{coef()}& extracts estimates&$\surd$\\
    \code{vcov()}& variance-covariance matrix&$\surd$\\
    \code{SE()}& standard error&$\surd$\\
    \alert{\code{summary()}}& shows summary of fitted model&\\
    \alert{\code{mse()}}& computes mean square error&\\
    \alert{\code{residuals()}}& extracts residuals&\\
    \alert{\code{fitted()}}& computes fitted values&\\
    \alert{\code{robweights()}}& robustness weights ($M$-estimators)&\\
    \alert{\code{scale()}}& estimate of scale ($M$-estimators)&\\
\end{tabular}
\vspace{1em}

{\small Note: $\surd$ indicates methods that are also available in the
    \code{survey} package.}
\end{frame}

% -----------------------------------------------------------------------------
\begin{frame}
\frametitle{What more?}
    \begin{itemize}
    \setlength\itemsep{0.5em}
    \item Implemented in the \code{C} language
%    \item Sums are computed as Kahan-type compensated sums
    \item \alert{$M$-estimators}
    \begin{itemize}
        \setlength\itemsep{0.25em}
        \item Huber and Tukey $\psi$-function
        \item Interface to add other $\psi$-functions: see
            \code{doc\_psifunction.html}
    \end{itemize}
%    \item Other functions like \code{weighted\_mad()}, \code{weighted\_IQR()}, etc.
    \item \alert{Vignettes}
    \begin{itemize}
        \setlength\itemsep{0.25em}
        \item Basic Robust Estimators
        \item Robust Horvitz--Thompson Estimator
    \end{itemize}
\end{itemize}
\end{frame}

% -----------------------------------------------------------------------------
\bgroup
\setbeamercolor{background canvas}{bg=fhnwblue}
\begin{frame}[c,plain]{}
\begin{center}
    \textcolor{white}{\large PART 2}\\
    \vspace{1em}
    \textcolor{white}{\Huge Robust Weighted Regression}
\end{center}
\end{frame}
\egroup

% -----------------------------------------------------------------------------
\begin{frame}[fragile]
\frametitle{Regression}
    \begin{itemize}
        \item Simple random sample of $n = 100$ counties in the U.S.
        \item Population: $N = 3\,141$ counties
        \item Data: Lohr (1999)
    \end{itemize}
    \begin{lstlisting}[style=in_big]
> head(counties[, c(2, 6, 7, 9, 10)], 2)
    \end{lstlisting}
    \vspace{-1em}
    \begin{lstlisting}[style=out_big]
    county  farmpop  numfarm  weights   fpc
1 Escambia      531      414    31.41  3141
2 Marshall     1592    15824    31.41  3141
    \end{lstlisting}
    where\\
    \vspace{0.5em}
    \begin{tabular}{ll|lll}
        \code{farmpop} & farm population && \code{weights} &weights\\
        \code{numfarm} & number of farms && \code{fpc} &population size\\
    \end{tabular}

\end{frame}

% -----------------------------------------------------------------------------
\begin{frame}[fragile]
\frametitle{Regression: Weighted least squares}
    % file_name <- "scatter_counties.pdf"
    % file_name_log <- "scatter_counties_log.pdf"
    % col_orange <- rgb(199, 90, 9, maxColorValue = 255)
    % library("robsurvey"); library("survey")
    % data(counties)
    % dn <- svydesign(ids = ~1, fpc = ~fpc, weights = ~weights,
    %     data = counties[counties$farmpop > 0, ])
    % pdf(file = file_name, height = 7, width = 5)
    % par(mar = c(5, 5, 4, 2))
    % plot(farmpop ~ numfarm, dn$variables, xlab = "numfarm", ylab = "farmpop",
    %     cex.axis = 1.4, cex.lab = 1.8)
    % dev.off()
    % system(paste("pdfcrop.exe", file_name, file_name))
    % pdf(file = file_name_log, height = 7, width = 5)
    % par(mar = c(5, 5, 4, 2))
    % plot(farmpop ~ numfarm, dn$variables, xlab = "numfarm (log)",
    %     ylab = "farmpop (log)", log = "xy", cex.axis = 1.4, cex.lab = 1.8)
    % points(farmpop ~ numfarm, dn$variables[3, ], pch = 19, col = col_orange)
    % arrows(10, 70, 3, 70, length = 0.15, col = col_orange)
    % dev.off()
    % system(paste("pdfcrop.exe", file_name_log, file_name_log))
    \begin{multicols}{2}
        \includegraphics[width=4.5cm]{figure/scatter_counties}
        \columnbreak
        \begin{itemize}
            \setlength\itemsep{1em}
            \item \alert{Model:} \code{farmpop} $\sim$ \code{numfarm}
            \item \alert{Variance:} (heteroscedasticity)
                \code{var = sqrt(numfarm)}
            \item \alert{Sampling design:} % i.e., "survey" must be loaded
\begin{verbatim}
dn <- svydesign(ids = ~1,
   fpc = ~fpc,
   weights = ~weights,
   data = subset(counties,
   numfarm > 0))
\end{verbatim}
            \item \alert{Weighted least squares}
        \end{itemize}
    \end{multicols}
\end{frame}

% -----------------------------------------------------------------------------
\begin{frame}[fragile]
    \frametitle{Regression: Weighted least squares {\small (cont'd)}}
    \begin{lstlisting}[style=in_small]
> svyreg(farmpop ~ numfarm, dn, var = ~sqrt(numfarm))
    \end{lstlisting}
    \vspace{-1em}
    \begin{lstlisting}[style=out_small]

Weighted least squares

Call:
svyreg(formula = farmpop ~ numfarm, design = dn,
       var = ~vi)

Coefficients:
(Intercept)      numfarm
    -53.998        1.839

Scale estimate: 99.25
    \end{lstlisting}
\end{frame}

% -----------------------------------------------------------------------------
\begin{frame}[fragile]
    \frametitle{Regression: Weighted least squares {\small (cont'd)}}
    \begin{itemize}
        \setlength\itemsep{1em}
        \item \alert{Methods}: \code{coef()}, \code{residuals()},
            \code{fitted()}, \code{plot()}, etc.
        \item \alert{Inference} under the model
            $y_i = \bm x_i^T \bm \theta + \sigma \sqrt{v_i} e_i$, $i \in U$
        \begin{itemize}
        {\normalsize
            \setlength\itemsep{1em}
            \vspace{0.5em}
            \item $\bm \theta$: super-population parameter
            \item $\bm \theta_N$: census parameter, finite-population parameter
            \item $\widehat{\bm \theta}_n$: sample-based estimator
        }
        \end{itemize}
        \item \alert{Design-based}: estimate $\bm \theta_N$
        \begin{verbatim}summary(..., mode = "design")\end{verbatim}
        \begin{verbatim}vcov(..., mode = "design")\end{verbatim}
    \end{itemize}
\end{frame}

% -----------------------------------------------------------------------------
\begin{frame}[fragile]
    \frametitle{Regression: Weighted least squares {\small (cont'd)}}
    \begin{itemize}
        \setlength\itemsep{1em}
        \item \alert{Model-based}: estimate $\bm \theta$ (ignore sampling
            design)\\
            %Let the LS estimator $\widehat{\bm \theta}_{LS}$ be an estimator
            %of $\bm \theta$\\
            %(conditional on sample $s$)\\
            \vspace{0.5em}
            \code{mode = "model"} in \code{summary()} and \code{vcov()}
        \item \alert{Compound design-model}: estimate $\bm \theta$\\
            %Let $\widehat{\bm \theta}_n$ be an estimator of $\bm \theta$\\
            \vspace{0.5em}
            \code{mode = "compound"} in \code{summary()} and \code{vcov()}
            \vspace{0.5em}
        % \begin{itemize}
        %     \setlength\itemsep{1em}
        %     {\normalsize
        %     \item Large sampling fraction $f$ (e.g., $f > 20\%$)
        %     \item Otherwise contribution of model-based variance is negligible
        %     }
        % \end{itemize}
    \end{itemize}
    \vspace{1em}
    {\small Literature: Rubin-Bleuer and Schriopu-Kratina (2005), Binder
        and Roberts (2009)}
 \end{frame}

% -----------------------------------------------------------------------------
\begin{frame}[fragile]
    \frametitle{Robust regression}
    \begin{multicols}{2}
        \includegraphics[width=4.5cm]{figure/scatter_counties_log}
        \columnbreak
        \begin{itemize}
            \setlength\itemsep{1em}
            \item \alert{Model:} \code{log(farmpop)} $\sim$ \code{log(numfarm)}
            \item \alert{Variance:} homoscedastic
            \item \alert{Outlier} $\rightarrow$ Robust regression
        \end{itemize}
    \end{multicols}
\end{frame}

% -----------------------------------------------------------------------------
\begin{frame}
\frametitle{Robust regression: Implementation}
\vspace{-0.5em}
\begin{center}
{ \footnotesize
\begin{tabular}{llrrrr}
    \multicolumn{2}{l}{Language} & $M$ & $GM$ & Reference\\
    \addlinespace
    \toprule
    \multicolumn{5}{l}{\alert{\textit{\footnotesize{with weighting}}}} \\
    \addlinespace
    \code{MASS}& \proglang{R} & $\surd$  & -- & Venables and Ripley (2002)\\
    \code{robust}& \proglang{R} & $\surd$ & -- & Wang et al. (2022)\\
    \code{robustbase}& \proglang{R} & $\surd$ & -- &
        M{\"a}chler et al. (2021)\\
    \code{robstat}& \proglang{Stata} & $\surd$ & -- & Jann et al. (2018)\\
    \code{robustreg}& \proglang{SAS}& $\surd$ & -- & SAS Institute (2022)\\
    \alert{\code{robsurvey}} & \alert{\proglang{R}} & \alert{$\surd$}
        & \alert{$\surd$} & \alert{this talk}\\
    \addlinespace
    \midrule
    \multicolumn{5}{l}{\textit{\footnotesize{without weighting}}} \\
    \code{robeth}& \proglang{R} & $\surd$ & $\surd$ & Marazzi (1993, 2020)\\
    \proglang{NAG}& \proglang{C} & $\surd$ & $\surd$ & NAG (2022)\\
    \code{GSL}& \proglang{C} & $\surd$ & -- & Galassi et al. (2019)\\
    \code{[?]regress} & \proglang{Stata} & $\surd$ & -- & Verardi and
        Croux (2009)\\
    \code{rreg}& \proglang{Stata}& $\surd$ & -- & StataCorp (2022)\\
    \code{robustfit}& \proglang{MATLAB} & $\surd$ & -- &The Math Works (2022)\\
\end{tabular}
}
\end{center}
\end{frame}

% -----------------------------------------------------------------------------
\begin{frame}[fragile]
    \frametitle{Robust regression: \emph{M}- and \emph{GM}-estimators}
     \alert{Function} \code{svyreg} followed by
    \begin{verbatim}
   _huberM(formula, design, k, var = NULL, ...)
  _huberGM(formula, design, k,
           type = c("Mallows", "Schweppe"),
           xwgt, var = NULL, ...)
    \end{verbatim}
    \vspace{-1.5em}
    \begin{itemize}
        \setlength\itemsep{0.5em}
        \item \code{k}: robustness tuning constant of Huber $\psi$-function
        \item \code{type}: Mallows or Schweppe \emph{GM}-estimator
        \item \code{xwgt}: downweight high-leverage observations
        \item \alert{Also} \code{_tukeyM()} and \code{_tukeyGM()} with Tukey
            $\psi$-function
    \end{itemize}
\end{frame}

% -----------------------------------------------------------------------------
% \begin{frame}[fragile]
%     \frametitle{Robust regression: \emph{M}- and \emph{GM}-estimators
%         {\small (cont'd)}}
%     \begin{lstlisting}[style=in_small]
% > svyreg_huberM(log(farmpop) ~ log(numfarm), dn,
%                 k = 2)
%     \end{lstlisting}
%     \vspace{-1.5em}
%     \begin{lstlisting}[style=out_small]

% Survey regression M-estimator (Huber psi, k = 2)

% Call:
% svyreg_huberM(formula = farmpop ~ numfarm,
%               design = dn, k = 2)

% IRWLS converged in 7 iterations

% Coefficients:
% (Intercept)      numfarm
%     -78.980        1.805

% Scale estimate: 388.4 (weighted MAD)
%     \end{lstlisting}
% \end{frame}

% -----------------------------------------------------------------------------
\begin{frame}
    \frametitle{Robust regression: \emph{M}- and \emph{GM}-estimators
        {\small (cont'd)}}
    \vspace{1em}
    \begin{itemize}
        \setlength\itemsep{1em}
        \item \alert{Utility methods:} \code{coef()}, \code{plot()}, etc.
        \item \alert{Inference:} \code{summary()} and \code{vcov()} (3 modes
            of inference)
        \item Vignette: \alert{Robust Weighted Regression}
    \end{itemize}
\end{frame}

% -----------------------------------------------------------------------------
\bgroup
\setbeamercolor{background canvas}{bg=fhnwblue}
\begin{frame}[c,plain]{}
\begin{center}
    \textcolor{white}{\large PART 3}\\
    \vspace{1em}
    \textcolor{white}{\Huge Summary and Outlook}
\end{center}
\end{frame}
\egroup

% -----------------------------------------------------------------------------
\begin{frame}
    \frametitle{Summary and outlook}
    \vspace{1em}
    \begin{itemize}
        \setlength\itemsep{0.75em}
        \item Package: \alert{52 exported functions} (+ 21 S3-methods)
        \vspace{0.5em}
        {\large
        \begin{itemize}
            \setlength\itemsep{0.5em}
            \item Robust generalized regression (GREG) estimator
            \item Tukey's weighted line
            \item ...
        \end{itemize}
        }
        \item Take away message: \alert{2 ``flavors''} of functions
        \vspace{0.5em}
        {\large
        \begin{itemize}
            \setlength\itemsep{0.5em}
            \item Bare-bone functions
            \item Survey methods (\code{survey} package required)
                $\Rightarrow$ variance
        \end{itemize}
        }
        \item \alert{What is missing? What methods do you need?}
    \end{itemize}
\end{frame}

% -----------------------------------------------------------------------------
\begin{frame}
    \frametitle{Summary and outlook {\small (cont'd)}}
    \vspace{1em}
    \begin{itemize}
        \setlength\itemsep{0.75em}
        \item Where can I find this slide deck?
        \vspace{0.5em}
        {\large
        \begin{itemize}
            \setlength\itemsep{0.5em}
            \item \alert{CRAN} webpage of \code{robsurvey}
            \item \alert{Link} to \alert{GitHub} \code{tobiasschoch/robsurvey}
            \item On GitHub: \alert{folder:}
                \code{/inst/doc/useR\_2022\_conference}
        \end{itemize}
        }
    \end{itemize}
\end{frame}

% -----------------------------------------------------------------------------
\bgroup
\setbeamercolor{background canvas}{bg=fhnwblue}
\begin{frame}[c,plain]{}
\begin{center}
    \textcolor{white}{\Huge I'm ready to take your questions!}
\end{center}
\end{frame}
\egroup

% -----------------------------------------------------------------------------
\begin{frame}
    \frametitle{Appendix: Datasets}
    {\footnotesize
    \begin{itemize}
        \setlength\itemsep{0.5em}
        \item \alert{counties}: Lohr SL (1999). \emph{Sampling: Design and
            Analysis}. Pacific Grove (CA): Duxbury Press, Appendix C.
        \item \alert{losdata}: Ruffieux C, Paccaud F, Marazzi A (2000).
            Comparing rules for truncating hospital length of stay.
            \emph{Casemix Quarterly} 2.
    \end{itemize}
    }
\end{frame}

% -----------------------------------------------------------------------------
\begin{frame}
\frametitle{Appendix: Software}
    {\footnotesize
    \begin{itemize}
        \setlength\itemsep{0.25em}
        \item \alert{GSL} Galassi M, Davies J, Theiler J, Gough B, Jungman G,
            Alken P, Booth M, Rossi F, Ulerich R (2019). GNU Scientific Library
            (release 2.6). 3rd edition.
        \item \alert{MASS} Venables WN and Ripley BD, 2002, Modern Applied
            Statistics with S. 4th edition, New York: Springer-Verlag.
        \item \alert{NAG} NAG (2022). The NAG Library for C. The Numerical
            Algorithms Group (NAG), Oxford. C library mark 28.3.
        \item \alert{[?]regress} Verardi V, Croux C (2009). Robust Regression
            in Stata. \emph{The Stata Journal} 9, 439--453.
            {\scriptsize \textcolor{orange_dark}{(Note: [?] is a wildcard for
            \code{m}, \code{s}, or \code{mm}; the methods thus read
            \code{mregress}, etc.)}}
        \item \alert{robeth} Marazzi A, 2020, robeth: R Functions for
            Robust Statistics. R package version 2.7-6.
        \item \alert{robstat} Jann B, Verardi V, Vermandele C, 2018. ROBSTAT:
            Stata module to compute robust univariate statistics. Statistical
            Software Components, Boston College Department of Economics.
    \end{itemize}
    }
\end{frame}

% -----------------------------------------------------------------------------
\begin{frame}
    \frametitle{Appendix: Software {\small (cont'd)}}
    {\footnotesize
    \begin{itemize}
        \setlength\itemsep{0.25em}
        \item \alert{robust} Wang J, Zamar R, Marazzi A, Yohai V,
            Salibián-Barrera M, Maronna R, Zivot E, Rocke D, Martin D,
            Mächler M, Konis K (2022). robust: A Port of the S-PLUS
            "Robust Library".  R package version 0.7-0.
        \item \alert{robustbase} Mächler M, Rousseeuw P, Croux C, Todorov V,
            Ruckstuhl A, Salibián-Barrera M, Verbeke T, Koller M, Conceicao ELT,
            Anna di Palma M (2021). robustbase: Basic Robust Statistics.
            R package version 0.93-9.
        \item \alert{robustfit} The Math Works, Inc. (2022). MATLAB. Version
            R2022a.
        \item \alert{robustreg} SAS Institute, Inc. (2020). SAS/STAT Software.
            SAS Institute Inc., Cary. Version 15.2.
        \item \alert{rreg} StataCorp (2022). Stata Statistical Software.
            StataCorp LLC, College Station. Release 17.
        \item \alert{survey} Lumley T (2021). survey: Analysis of Complex
            Survey Samples. R package version 4.1-1.
    \end{itemize}
    }
\end{frame}

% -----------------------------------------------------------------------------
\begin{frame}
    \frametitle{Appendix: Literature}
    {\footnotesize
    \begin{itemize}
        \setlength\itemsep{0.25em}
        \item \alert{Beaumont JF and Rivest LP (2009)}. Dealing with outliers
            in survey data, in \emph{Sample Surveys: Theory, Methods and
            Inference}, ed. by Pfeffermann D and Rao CR, Amsterdam: Elsevier,
            vol. 29A of Handbook of Statistics, chap. 11, 247–280.
        \item \alert{Binder DA and Roberts G (2009)}. Design- and Model-Based
            Inference for Model Parameters. In: \emph{Sample Surveys: Inference
            and Analysis ed. by Pfeffermann}, D. and Rao, C. R. Volume 29B
            of Handbook of Statistics, Amsterdam: Elsevier, Chap. 24, 33--54.
        \item \alert{Hampel FR, Ronchetti EM, Rousseeuw PJ and Stahel WA
            (1986)}.  \emph{Robust Statistics: The Approach Based on Influence
            Functions}.  New York: John Wiley and Sons.
        \item \alert{Hulliger B (1995)}. Outlier Robust Horvitz–Thompson
            Estimators.  \emph{Survey Methodology} 21, 79--87.
        \item \alert{Lee H (1995)}. Outliers in Business Surveys, in
            \emph{Business Survey Methods}, ed. by Cox BG,
            Binder DA, Chinnappa BN, Christianson A, Colledge MJ, Kott PS,
            New York: John Wiley and Sons, chap. 26, 503--526.
    \end{itemize}
    }
\end{frame}

% -----------------------------------------------------------------------------
\begin{frame}
    \frametitle{Appendix: Literature {\small (cont'd)}}
    {\footnotesize
    \begin{itemize}
        \setlength\itemsep{0.25em}
        \item \alert{Lumley T (2010)}. \emph{Complex Surveys: A Guide to
            Analysis Using R.}, Hoboken (NJ): John Wiley and Sons.
        \item \alert{Rubin-Bleuer S and Schiopu-Kratina I (2005)}. On the
            Two-phase framework for joint model and design-based inference.
            \emph{The Annals of Statistics} 33, 2789--2810.
    \end{itemize}
    }
\end{frame}


% -----------------------------------------------------------------------------
\end{document}
